\section{Related work}
\label{sec:related}

Per-paper narration of the works cited here is in the supplementary literature review; this section keeps what bears
on a novelty claim.

\subsection{Classifier ensemble weighting}
\label{sec:rel:weighting}

Weighted voting and weighted probability averaging are long established
\citep{hansen1990neural,kittler1998combining,dietterich2000ensemble,kuncheva2014combining}; the weights come from
stacked generalization
\citep{wolpert1992stacked,breiman1996stacked,leblanc1996combining,vanderlaan2007super,ledell2016auc}, greedy forward
selection \citep{caruana2004ensemble}, numerical optimization
\citep{zhang2011sparse,yao2018stacking,shahhosseini2022optimizing,hassan2026weighted,xu2026pseo} or theory
\citep{fumera2005theoretical}. Throughout, the simplex is almost always respected as a constraint rather than designed
over as an experimental region, and essentially all of this work is single-objective.

\subsection{Multiobjective ensemble optimization}
\label{sec:rel:moo}

Multiobjective ensemble learning divides by decision variable: the \emph{members}
\citep{abbass2003pareto,chandra2006ensemble,jin2008pareto}; a \emph{subset}
\citep{dossantos2008dynamic,qian2015pareto}; or --- the family relevant here --- \emph{continuous combination
weights} under conflicting criteria, optimized by metaheuristics
\citep{ekbal2011multiobjective,ekbal2013combining,bandyopadhyay2008amosa,onan2016multiobjective,rezaei2019multi,zhao2018multiobjective,vrugt2006multi,galvan2026simultaneous}
and surveyed in \citet{ribeiro2020ensemble}.

This family is therefore \emph{not} new, and we do not claim it. Absent from this cluster is any designed experiment
and fitted response surface over the weight simplex, any classical scalarization method such as NBI, and --- with the
partial exception of the sparsity term in \citet{zhao2018multiobjective} --- any treatment of deployment cost as an
objective. Fronts are rarely replicated over resampled partitions with corrected inference
\citep{demsar2006statistical}.

\subsection{Mixture experiments and response surfaces}
\label{sec:rel:mixture}

Mixture experiments address exactly the constraint set of \cref{eq:simplex}, through simplex-lattice designs and
Scheffé's canonical polynomials
\citep{scheffe1958mixtures,scheffe1963simplexcentroid,cornell2002mixtures,myers2016rsm}; because $\sum_i w_i = 1$,
their coefficients are not component effects in the ordinary sense \citep{piepel1982component}.

The classical reading of a positive interaction coefficient is \emph{synergistic blending}
\citep{cornell2002mixtures}: a statement relative to the chord, not relative to the better pure component, and the
distinction turns out to matter (\cref{sec:results:coefficients}).

Mixture designs are beginning to appear in machine learning
\citep{karl2023svem,leal2025mixtureweights,xie2023doremi,ye2025mixinglaws,liu2025regmix,chen2025aioli,thudi2025mixmin},
but share our formalism and not our object: none models the performance of a weighted \emph{prediction} ensemble.

The exception, and the closest prior art to this study, is \citet{kwon2024ensemble}: the weights of five heterogeneous
classifiers as mixture components, a Scheffé canonical polynomial of accuracy or $F_1$ fitted over a mixture design,
reduced by backward elimination and maximized subject to $\sum_i w_i = 1$, on eleven tabular binary datasets. The
simplex formulation, the mixture design over classifier weights and the Scheffé surrogate are therefore established
prior art for classifiers, and we claim none of them. What that study does not do is what this paper is about: it
optimizes one scalar metric at a time, so it has no Pareto front, no anchors and no NBI; it validates the fitted
polynomial only in-sample, against no separately sampled reference front; and it carries no cost objective.

\subsection{The DoE--RSM--NBI framework and its application to ensemble weights}
\label{sec:rel:lineage}

Normal Boundary Intersection \citep{das1998nbi} was introduced to obtain a distributed representation of the Pareto
boundary, by solving for each convex weight vector a subproblem that advances from the convex hull of individual
minima along a quasi-normal direction. It
is widely combined with response-surface metamodels in engineering optimization, and a substantial body of work applies
NBI to polynomial surfaces fitted from designed experiments
\citep{brito2014nbimse,costa2016nbimmse,lopes2016rpdmnbi,naves2017nbirfs,belinato2019mnbipca,almeida2022taguchinbi}.

Within that tradition, mixture designs have been used to schedule the NBI weights and to model post-Pareto criteria
over the weight simplex. \citet{mendes2016portfolio_mde} and \citet{leal2022portfolio_doptimal_mde} model portfolio
proportions with mixture designs; \citet{bacci2019nbi_mixture} builds a simplex lattice over time-series
forecast-combination weights, fits Scheffé models to factor scores of residual accuracy metrics and applies NBI to
those models; \citet{paula2019gatuning} crosses a mixture design over objective weights with genetic-algorithm
hyperparameters. \citet{pereira2025hybrid}, of which the present first author is a co-author, extends the framework
with a mixture-design post-Pareto stage in which Scheffé polynomials of a generalized distance and of the entropy of
the weights are themselves optimized by a second NBI. Its anchors, like those of the rest of this literature, are the
individual optima of the surrogate surfaces --- the configuration we label NBI-A. The first author's master's
dissertation \citep{ribeiro2026dissertation} applies the same framework to hyperparameter optimization.

Most directly, \citet{moreira2021ann_ensemble_mde} determine neural-network ensemble weights by a mixture design, and
\citet{rocha2025ensemblestlf} combine a $\{3,5\}$ simplex-lattice design over the mixing weights of three neural
networks, a factor-analytic surrogate of correlated error metrics, NBI over those surrogates, and an entropy-based
post-Pareto choice, for short-term load forecasting. That is the present pipeline applied to a regression ensemble, and
its stated first contribution is casting ensemble-weight definition as a structured mixture-design problem.

\textbf{We therefore claim no part of this construction as new.} Our contribution begins where this literature stops:
none of these studies validates its mixture surrogate on unseen compositions, computes anchors from anything other than
the surrogate, revalidates its Pareto candidates on the true objectives, scores them against a separately sampled
reference, or replicates over resampled partitions.

\subsection{Surrogate-assisted multiobjective optimization}
\label{sec:rel:surrogate}

The broader surrogate-assisted literature \citep{jin2011surrogate,tabatabaei2015survey,chugh2019survey,deb2019taxonomy}
answers the risk that a surrogate misleads with \emph{model management}, and that answer is \emph{iterative}: ParEGO,
MOEA/D-EGO, SMS-EGO and K-RVEA \citep{knowles2006parego,zhang2010moeadego,ponweiser2008smsego,chugh2018krvea} refit on
true-objective infill inside the search loop, so a false surrogate optimum is corrected later. The pipeline studied
here is that taxonomy's other member: the \emph{sequential} framework \citep{tabatabaei2015survey} --- fit the
surrogate once from a fixed design, then optimize it, with no infill --- identified there as most exposed to surrogate
error. That is what the lineage of
\cref{sec:rel:lineage} does and what we replicate, so our diagnosis is a diagnosis of \emph{one-shot, no-infill}
surrogate pipelines. It is not a test of model-managed, infill-based surrogate-assisted
optimizers and must not be read as one: by the standards of that literature the exposure of the surrogate-anchored arm
is anticipated rather than surprising, and our real-anchor arm is best understood as a single targeted correction ---
one model-management step applied to the anchors alone --- rather than as adaptive infill.

Extreme-point reliability has been studied, but on the other side of the surrogate divide.
\citet{isermann1988payoff} showed computationally that payoff-table minima built from individual optima can differ
substantially from the true minimum criterion values over the efficient set; nadir and front-range mis-estimation
degrades performance across MOEAs \citep{deb2010nadir,he2021normalization}; and \citet{herrmann2026nonextreme} treat
the \emph{choice} of individual minima as a design decision, replacing extreme with non-extreme minima to refine the
utopia--nadir hyperbox.
That misplaced anchors damage a scalarization is therefore established, and we claim neither
the mechanism nor its direction as a discovery. In all of that work the error arises from degeneracy or search
difficulty on the \emph{true} objectives; a metamodel is never the culprit. The normalized normal constraint
\citep{messac2003normalized}, introduced partly to make anchor-based normalization better behaved, is the
textbook remedy for the pathology we report; our protocol was frozen with NBI as the front generator, so we do not
evaluate it here and record it as the most direct methodological follow-up.

What we did not find is a study computing a CHIM method's anchors both ways on one problem and comparing the
resulting fronts --- the surrogate \emph{provenance} of the anchor error, and the controlled contrast, rather than the
fact that anchors matter. The closest external analogue we located is the NBI-A pattern itself:
\citet{abolghasemian2022haulage} runs a modified NBI on regression metamodels validated by fit statistics alone,
its Pareto solutions never re-evaluated on the simulation (supplementary literature review). At the other
extreme, \citet{gellerich2023doenbi} derive the experimental plan from the NBI subproblems themselves and evaluate each
subproblem by physical test, with no metamodel anywhere and anchors that are real experimental optima. A metamodel-free
NBI is thus published prior art and the idea of our NBI-C arm is not claimed as new; because that study has no
surrogate counterpart on its problem, it makes no comparison. The two papers bracket our NBI-A and NBI-C arms without
either of them running the contrast.

\subsection{Cost-aware ensemble deployment}
\label{sec:rel:cost}

Prediction-time cost has been an explicit objective since \citet{turney2000types}, priced as a fixed \emph{deployed
set} \citep{zhou2002ensembling,ji2023pruning,maier2024hardware,maier2026hapens,gunasekaran2022cocktail,moreira2026energy},
as \emph{per-example adaptive execution}
\citep{xu2013cstc,trapeznikov2013multistage,nan2017adaptive,chen2020frugalml}, or as a \emph{continuous penalty} in a
scalarized training objective \citep{xu2012greedy,zhang2011sparse}; automated machine-learning systems expose
inference-time constraints during ensemble construction
\citep{erickson2020autogluon,feurer2022autosklearn,borchert2022paretoselect}. For post-hoc ensembles over cached
tabular predictions --- our own setting --- \citet{maier2024hardware} already return a Pareto front of accurate and
efficient ensembles with inference time in the objective vector over 83 tabular classification datasets, and
\citet{maier2026hapens} extend that line, name the objective \emph{deployment cost} and report memory usage as their
most effective cost metric. Inference cost as an objective of a post-hoc weighted ensemble is therefore prior art, and
we claim no part of it.

The distinction we need falls between these categories. That the accounting itself can reorder a comparison is
already known: \citet{wang2022committees} place two cost accountings side by side --- every member executed, versus
early exit --- and show that which committee looks best changes with the accounting. Their contrast, however, is
between ensemble and cascade \emph{architectures}, not between two cost functions applied to one and the same weight
vector. Narrowed to that resolution, to the best of our literature search no prior work contrasts the step cost a
deployed support pays with the continuous relaxation the optimizers minimize \emph{for one weight vector of one
weighted soft-voting ensemble}, and reports that the two disagree about which method wins. Where a related question
has been settled it went the other way: for budgeted learning with feature re-use, the support-cost programme is
provably solved by its linear relaxation \citep{nan2015budgeted}. Our setting does not
inherit that guarantee, and \cref{sec:results:cost} shows the two definitions disagreeing.

\subsection{The gap}
\label{sec:rel:gap}

Ensemble weights on a simplex, mixture designs over them, and Scheffé surrogates of ensemble performance all exist,
and for classifier ensembles specifically they exist together in \citet{kwon2024ensemble}; adding NBI on those
surrogates exists in \citet{rocha2025ensemblestlf} for a regression ensemble and in \citet{bacci2019nbi_mixture} for
forecast combination. Multiobjective optimization of continuous ensemble weights exists. Cost-aware ensemble selection
with deployment cost as a Pareto objective exists \citep{maier2024hardware,maier2026hapens}. Metamodel-free NBI on
measured objectives exists \citep{gellerich2023doenbi}, and that anchor misplacement degrades a scalarization is
established \citep{isermann1988payoff,deb2010nadir,herrmann2026nonextreme}. What does not exist, to the best of our
literature search, is an evaluation of whether this pipeline can be trusted once it is both multiobjective and applied
to classifiers: no external validation of the mixture surrogate on unseen compositions, no computation of the anchors
both ways on one problem so that surrogate-derived and real anchors can be contrasted, no metamodel-free arm run
alongside the surrogate arms as a control, no revalidation of candidates on the true objectives against an
empirical reference with a method-independent core, no replicated partition study of the whole pipeline, and no
contrast, for one and the same weight vector, between the cost that is optimized and the cost that is paid.

The gap is therefore not a missing method. It is the absence of the controls that would tell a practitioner when the
existing method's output can be believed --- which is what this paper supplies, for the one-shot, no-infill form of
the pipeline described above.
